\section*{Your turn}

This chapter has two tags, and the exercises use both. Start the database once
and leave it running:

\begin{minted}{text}
git clone https://github.com/Giang-Dang/mosaic-graph
cd mosaic-graph
docker compose up -d mosaic-db
git checkout ch04
pwsh scripts/verify.ps1 -KeepDatabase
\end{minted}

The script now brings PostgreSQL up itself. It asserts three lookups and three
statements where chapter~\ref{ch:the-life-of-a-request}'s asserted a hundred and
forty-six lookups and knew nothing about statements, because Mosaic had no
database then. At tag \code{ch04-ef} the same script asserts a hundred and
forty-six of both. Run the service with its console visible for all of these.

\begin{enumerate}
  \item \textbf{Watch the number change under you.} Send the
  catalog-with-reviews query at tag \code{ch04} and read the timeline. Now
  \code{git checkout ch04-ef} and send the same query. Two things moved and one
  did not. Write down which, and say in one sentence why the one that did not
  move is the interesting one.

  \item \textbf{Find the split.} Warm the service at tag \code{ch04}, then send
  the query four hundred times and count what the timeline reported:

  \begin{minted}[fontsize=\footnotesize]{text}
grep -oE "[0-9]+ SQL\)" log | sort -n | uniq -c
  \end{minted}

  You are looking for the runs that report four. Turn the database command log
  up to \code{Information} and work out which statement doubled. Then explain,
  from section~\ref{sec:ch04-when-a-batch-goes}, why it is that one and not one
  of the others.

  \item \textbf{Break the batching deliberately.} In
  \code{ReviewDataLoaders.cs}, change the reviews loader so that it returns a
  dictionary of arrays rather than a lookup, and make the service return one. It
  compiles, and it still batches. Send the catalog-with-reviews query and count
  the errors. Which three products, and why those three?

  \item \textbf{Take the identifier away.} On \code{ProductNode.GetId}, delete
  the argument to the \code{Parent} attribute so that it no longer names the
  identifier, then ask \code{browseProducts} for its nodes' \code{id} and
  \code{title}. Before you look at the response, predict the HTTP status. Then
  run the Postman collection and watch which single assertion catches it.

  \item \textbf{Cost the count.} Turn the database command log up and send a
  page of three with and without \code{totalCount}. Count the statements. Now
  remove the connection attribute from the field, export the schema, and check
  whether \code{totalCount} left the SDL. The answer to that last question is
  the point of the exercise.

  \item \textbf{Ask for something the projection cannot reach.} Send
  \code{browseProducts} asking only for \code{nodes \{ price \{ amount \} \}}
  and read the statements. How many, and what is in the first one? Then work out
  why \code{price} could not have been a column in it even in principle.

  \item \textbf{Give a lookup a real cost.} Mosaic's PostgreSQL is on your own
  machine, which is the friendliest network a query will ever see. Starve the
  container of processor time with \code{docker update --cpus 0.25 mosaic-db}
  and re-run exercise 1 on both tags. The ratio between them is the number this
  chapter claims travels; find out whether it does, and put the container back
  afterwards.
\end{enumerate}

Exercise 4 is the one worth doing even if you skip the rest. Everything else
here fails loudly: a wrong return type produces errors, and a missing
configuration produces HC0018. A projection that drops a key produces a
well-formed response with the right
number of elements in it, a 200, and no errors key, and it does so only when a
client asks for the field rather than when your smoke test does. From
chapter~\ref{ch:enter-the-router} onward there is another process between the
client and this data with its own opportunities to lose a field quietly, and the
habit that survives all of them is the same one: assert on the value, not on the
shape.
